\documentclass{amsart}
% For dealing with references we use the comment environment
\usepackage{verbatim}
\newenvironment{reference}{\comment}{\endcomment}
%\newenvironment{reference}{}{}
\newenvironment{slogan}{\comment}{\endcomment}
\newenvironment{history}{\comment}{\endcomment}

% For commutative diagrams we use Xy-pic
\usepackage[all]{xy}

% We use 2cell for 2-commutative diagrams.
\xyoption{2cell}
\UseAllTwocells

% We use multicol for the list of chapters between chapters
\usepackage{multicol}

% This is generally recommended for better output
\usepackage{lmodern}
\usepackage[T1]{fontenc}

% For cross-file-references

% Package for hypertext links:
\usepackage{hyperref}

% For any local file, say "hello.tex" you want to link to please

% Theorem environments.
%
\theoremstyle{plain}
\newtheorem{theorem}[subsection]{Theorem}
\newtheorem{proposition}[subsection]{Proposition}
\newtheorem{lemma}[subsection]{Lemma}

\theoremstyle{definition}
\newtheorem{definition}[subsection]{Definition}
\newtheorem{example}[subsection]{Example}
\newtheorem{exercise}[subsection]{Exercise}
\newtheorem{situation}[subsection]{Situation}

\theoremstyle{remark}
\newtheorem{remark}[subsection]{Remark}
\newtheorem{remarks}[subsection]{Remarks}

\numberwithin{equation}{subsection}

% Macros
%
\def\lim{\mathop{\mathrm{lim}}\nolimits}
\def\colim{\mathop{\mathrm{colim}}\nolimits}
\def\Spec{\mathop{\mathrm{Spec}}}
\def\Hom{\mathop{\mathrm{Hom}}\nolimits}
\def\Ext{\mathop{\mathrm{Ext}}\nolimits}
\def\SheafHom{\mathop{\mathcal{H}\!\mathit{om}}\nolimits}
\def\SheafExt{\mathop{\mathcal{E}\!\mathit{xt}}\nolimits}
\def\Sch{\mathit{Sch}}
\def\Mor{\mathop{\mathrm{Mor}}\nolimits}
\def\Ob{\mathop{\mathrm{Ob}}\nolimits}
\def\Sh{\mathop{\mathit{Sh}}\nolimits}
\def\NL{\mathop{N\!L}\nolimits}
\def\CH{\mathop{\mathrm{CH}}\nolimits}
\def\proetale{{pro\text{-}\acute{e}tale}}
\def\etale{{\acute{e}tale}}
\def\QCoh{\mathit{QCoh}}
\def\Ker{\mathop{\mathrm{Ker}}}
\def\Im{\mathop{\mathrm{Im}}}
\def\Coker{\mathop{\mathrm{Coker}}}
\def\Coim{\mathop{\mathrm{Coim}}}

% Boxtimes
%
\DeclareMathSymbol{\boxtimes}{\mathbin}{AMSa}{"02}

%
% Macros for moduli stacks/spaces
%
\def\QCohstack{\mathcal{QC}\!\mathit{oh}}
\def\Cohstack{\mathcal{C}\!\mathit{oh}}
\def\Spacesstack{\mathcal{S}\!\mathit{paces}}
\def\Quotfunctor{\mathrm{Quot}}
\def\Hilbfunctor{\mathrm{Hilb}}
\def\Curvesstack{\mathcal{C}\!\mathit{urves}}
\def\Polarizedstack{\mathcal{P}\!\mathit{olarized}}
\def\Complexesstack{\mathcal{C}\!\mathit{omplexes}}
% \Pic is the operator that assigns to X its picard group, usage \Pic(X)
% \Picardstack_{X/B} denotes the Picard stack of X over B
% \Picardfunctor_{X/B} denotes the Picard functor of X over B
\def\Pic{\mathop{\mathrm{Pic}}\nolimits}
\def\Picardstack{\mathcal{P}\!\mathit{ic}}
\def\Picardfunctor{\mathrm{Pic}}
\def\Deformationcategory{\mathcal{D}\!\mathit{ef}}

\setlength{\emergencystretch}{3em}
\begin{document}
\title[Stacks Project, Tag 0B5V]{370 --- Ampleness descends through finite surjective maps of proper Noetherian schemes}
\maketitle
\noindent Copyright (C) 2005--2025 Johan de Jong. Stacks Project authors. Source: \href{https://stacks.math.columbia.edu/tag/0B5V}{Tag 0B5V}.

\noindent Editorial difficulty note (not author text). Difficulty H2: The difficult implication cannot assume the finite surjection is flat or use section norms directly. The proof builds a cohomological property of twisted coherent sheaves and verifies generic-support devissage by finite pushforward, projection and higher-cohomology vanishing on the covering scheme..

\noindent Editorial provenance note (not author text). History: excerpt from revision \texttt{a0444\allowbreak 6e57e\allowbreak c1fbc\allowbreak 252a8\allowbreak 71afc\allowbreak ec775\allowbreak 2fb28\allowbreak 07b14}, coherent.tex, lines 4547--4617. Reference presentation was adapted; original-author context and core support blocks were retained or added. The mathematical wording is unchanged. Licensed under GNU FDL 1.2 or later; no invariant sections or cover texts. See ../licenses/COPYING. Context blocks reproduce author definitions and supporting results; they are not additional counted problems.

\begin{lemma}
\label{lemma-surjective-finite-morphism-ample}
Let $R$ be a Noetherian ring. Let $f : Y \to X$ be a morphism of
schemes proper over $R$. Let $\mathcal{L}$ be an
invertible $\mathcal{O}_X$-module. Assume $f$ is finite and surjective.
Then $\mathcal{L}$ is ample if and only if $f^*\mathcal{L}$ is ample.
\end{lemma}

\begin{proof}
The pullback of an ample invertible sheaf by a quasi-affine morphism
is ample, see Morphisms, Lemma
\href{https://stacks.math.columbia.edu/tag/0892}{lemma pullback ample tensor relatively ample (Tag 0892)}.
This proves one of the implications as a finite morphism is affine
by definition.

\medskip\noindent
Assume that $f^*\mathcal{L}$ is ample. Let $P$ be the following property on
coherent $\mathcal{O}_X$-modules $\mathcal{F}$: there exists an $n_0$
such that $H^p(X, \mathcal{F} \otimes \mathcal{L}^{\otimes n}) = 0$
for all $n \geq n_0$ and $p > 0$. We will prove that $P$ holds
for any coherent $\mathcal{O}_X$-module $\mathcal{F}$, which implies
$\mathcal{L}$ is ample by Lemma \href{https://stacks.math.columbia.edu/tag/0B5U}{lemma vanshing gives ample (Tag 0B5U)}.
We are going to apply Lemma \href{https://stacks.math.columbia.edu/tag/01YM}{lemma property higher rank cohomological (Tag 01YM)}.
Thus we have to verify (1), (2) and (3) of that lemma for $P$.
Property (1) follows from the long exact cohomology sequence associated
to a short exact sequence of sheaves and the fact that tensoring with
an invertible sheaf is an exact functor. Property (2) follows since
$H^p(X, -)$ is an additive functor. To see (3) let $Z \subset X$ be
an integral closed subscheme with generic point $\xi$.
Let $\mathcal{F}$ be a coherent sheaf on $Y$ such that
the support of $f_*\mathcal{F}$ is equal to $Z$
and $(f_*\mathcal{F})_\xi$ is annihilated by $\mathfrak m_\xi$,
see Lemma \href{https://stacks.math.columbia.edu/tag/01YO}{lemma finite morphism Noetherian (Tag 01YO)}. We claim that
taking $\mathcal{G} = f_*\mathcal{F}$ works. We only have to verify
part (3)(c) of Lemma \href{https://stacks.math.columbia.edu/tag/01YM}{lemma property higher rank cohomological (Tag 01YM)}.
Hence assume that $\mathcal{J} \subset \mathcal{O}_X$ is a
quasi-coherent sheaf of ideals such that
$\mathcal{J}_\xi = \mathcal{O}_{X, \xi}$.
A finite morphism is affine hence by
Lemma \href{https://stacks.math.columbia.edu/tag/01YP}{lemma affine morphism projection ideal (Tag 01YP)} we see that
$\mathcal{J}\mathcal{G} = f_*(f^{-1}\mathcal{J}\mathcal{F})$.
Also, as pointed out in the proof of
Lemma \href{https://stacks.math.columbia.edu/tag/01YP}{lemma affine morphism projection ideal (Tag 01YP)} the sheaf
$f^{-1}\mathcal{J}\mathcal{F}$ is a coherent $\mathcal{O}_Y$-module.
As $\mathcal{L}$ is ample we see from Lemma \href{https://stacks.math.columbia.edu/tag/0B5U}{lemma vanshing gives ample (Tag 0B5U)}
that there exists an $n_0$ such that
$$
H^p(Y, f^{-1}\mathcal{J}\mathcal{F}
\otimes_{\mathcal{O}_Y} f^*\mathcal{L}^{\otimes n}) = 0,
$$
for $n \geq n_0$ and $p > 0$. Since $f$ is finite, hence affine, we see that
\begin{align*}
H^p(X, \mathcal{J}\mathcal{G} \otimes_{\mathcal{O}_X}
\mathcal{L}^{\otimes n})
& =
H^p(X, f_*(f^{-1}\mathcal{J}\mathcal{F}) \otimes_{\mathcal{O}_X}
\mathcal{L}^{\otimes n}) \\
& =
H^p(X, f_*(f^{-1}\mathcal{J}\mathcal{F} \otimes_{\mathcal{O}_Y}
f^*\mathcal{L}^{\otimes n})) \\
& =
H^p(Y, f^{-1}\mathcal{J}\mathcal{F} \otimes_{\mathcal{O}_Y}
f^*\mathcal{L}^{\otimes n}) = 0
\end{align*}
Here we have used the projection formula
(Cohomology, Lemma \href{https://stacks.math.columbia.edu/tag/01E8}{lemma projection formula (Tag 01E8)}) and
Lemma \href{https://stacks.math.columbia.edu/tag/089W}{lemma relative affine cohomology (Tag 089W)}.
Hence the quasi-coherent subsheaf $\mathcal{G}' = \mathcal{J}\mathcal{G}$
satisfies $P$. This verifies property (3)(c) of
Lemma \href{https://stacks.math.columbia.edu/tag/01YM}{lemma property higher rank cohomological (Tag 01YM)} as desired.
\end{proof}
\end{document}
